\section*{Reischl Meeting Notizen}


\newpage

\section*{Notizen:}

\section*{Modellanalyse}

\begin{enumerate}
    \item \textbf{Finanzdaten bereitstellen:}
          \begin{itemize}
              \item Kurse und eventuell weitere Kennzahlen sammeln
              \item Entscheidung über Darstellungsform: vermutlich Tagesabschlüsse
              \item Trennung von Trainings- und Testdaten
          \end{itemize}

    \item \textbf{Interfaces:}
          \begin{itemize}
              \item Kursvorhersage (Strategie für Kauf und Verkauf implementieren)
              \item oder direkte Kauf-/Verkaufsentscheidungen treffen
          \end{itemize}

    \item \textbf{Interface implementieren:}
          \begin{itemize}
              \item Ermöglicht das Testen verschiedener Modelle
          \end{itemize}
    \item \textbf{Benchmark bauen und Bewertungskriterien festlegen}
    \item \textbf{Aktuelle Modelle testen und bewerten}
    \item \textbf{Modelle auf Aktien und Indizes testen, vielleicht auch mal auf Gold}
    \item \textbf{Testen auf Tagesabschlüssen, wenn möglich auf Minuten/Sekunden Daten}
\end{enumerate}

\section*{Modellentwicklung}

\textbf{Ideen:}

\begin{enumerate}
    \item \textbf{Matrix zur Unternehmenskorrelation:}
          \begin{itemize}
              \item Beispiel: Steigt Nvidia, sinkt eventuell AMD $\Rightarrow$ Korrelation $=-0.2$
          \end{itemize}

    \item \textbf{Zeitreihenanalyse:}
          \begin{enumerate}
              \item Umwandlung in stationäre Daten (\textit{stationary analysis})
              \item Nutzung einfacher Attention-Mechanismen auf alten Daten $\Rightarrow$ Gewichtungsmatrix (ARIMA)
              \item Einschätzung der Schwankung mit GARCH
              \item Eventuell Rauschfiltern, Saisonalität, Intra-Day-Zyklen finden durch Fourier(?)
          \end{enumerate}

    \item \textbf{Ziel:}
          \begin{itemize}
              \item Kauf-/Verkaufsentscheidungen
              \item Kursvorhersage
              \item oder tatsächliche Unternehmenswertschätzung
              \item $\Rightarrow$ Aus all diesem Ableitung der Kaufentscheidung
          \end{itemize}
\end{enumerate}

\section*{Klassifikation und Entscheidungsfindung}

Vielleicht Einsatz einer \textbf{Support Vector Machine (SVM)} zur Trennung verschiedener Features für die Klassifikation und Kaufentscheidung.
(Im Endeffekt muss eine Kaufentscheidung getroffen werden, manche Modelle liefern aber Kursvorhersage oder Unternehmenswert. Daraus muss trinäre Kauf/Verkauf/Halten-Entscheidung getroffen werden)

\begin{itemize}
    \item Jede Dimension ist ein Trainingselement: (ARIMA, GARCH, Tweets, Unternehmenswert)
    \item SVM trennt in \textit{Kaufen} und \textit{Nicht Kaufen}
    \item falls SVM mit 3 Teilungen möglich, dann direkt in Kaufen, Halten, verkaufen teilen, ansonsten in Kaufen, Nicht-Kaufen und dann Nicht-Kaufen nochmal in Halten, Verkaufen teilen.
    \item Daraus resultieren konkrete Kaufentscheidungen
    \item \textbf{ToDo:} Risikomanagement analytisch bestimmen und Positionsgröße berechnen um Handelsentscheidungen zu konkretisieren
\end{itemize}

\section*{Nachrichtenanalyse zur Kaufentscheidungsfindung}

\begin{itemize}
    \item Verwendung von sozialen Medien zur Entscheidung ob Kurs fällt oder steigt
\end{itemize}

\subsection*{Indivdual}
\begin{itemize}
    \item Einzelne Accounts scrapen (Elon Musk, Trump), Modell entwerfen, wie dieser Tweet den Kurs welches Unternehmens verändern wird, eventuell mit Transformer, LLM, am liebsten mit simplerem Model
    \item eventuell Sentiment-Analysis, oder z.B. ähnlich wie Sentiment-Analysis für Aktienkurse selbst schreiben.
    \item Analyse dann in Kaufentscheidung einfließen lassen
\end{itemize}

\subsection*{Gemeinschaft}
\begin{itemize}
    \item Gesamte Kanäle, Twitter Trends, Wallstreetbets, finanzen.net o.Ä. scrapen
    \item daraus quantitative Analyse, welchen Trend welche Aktien grade nehmen.
    \item Sentiment dafür gute Wahl, da Stock bekannt und nur Trend wichtig (einzelne Ausreiser vernachlässigbar)
\end{itemize}